% Report: Bermudan swaptions in a one-factor HJM model.
% Build: pdflatex report.tex (twice). Figures read their data from data/*.dat, produced by the page's JavaScript engine.
\documentclass[11pt,a4paper]{article}
\usepackage[a4paper,margin=2.4cm]{geometry}
\usepackage[T1]{fontenc}
\usepackage[utf8]{inputenc}
\usepackage{lmodern}
\usepackage{microtype}
\usepackage{amsmath,amssymb,amsthm,mathtools,bm}
\usepackage{booktabs}
\usepackage{enumitem}
\usepackage{xcolor}
\usepackage{pgfplots}
\pgfplotsset{compat=1.17}
\usepackage{pgfplotstable}
\usepackage{caption}
\usepackage{algorithm}
\usepackage{algpseudocode}
\usepackage[colorlinks=true,linkcolor=blue!50!black,citecolor=blue!50!black,urlcolor=blue!50!black]{hyperref}
\usepackage{cleveref}
\hypersetup{pdftitle={Bermudan swaptions in a one-factor HJM model},pdfauthor={Mohammed-Yassine Barnicha}}

\newtheorem{proposition}{Proposition}
\newtheorem{lemma}[proposition]{Lemma}
\theoremstyle{definition}
\newtheorem{definition}{Definition}
\theoremstyle{remark}
\newtheorem{remark}{Remark}

\newcommand{\E}{\mathbb{E}}
\newcommand{\Q}{\mathbb{Q}}
\newcommand{\N}{\mathcal{N}}
\newcommand{\Bb}{\bar B}
\newcommand{\dd}{\,\mathrm{d}}
\pgfplotsset{every axis/.append style={width=0.48\textwidth,height=5.2cm,grid=major,grid style={gray!25},
  scaled ticks=false,xticklabel style={/pgf/number format/fixed},yticklabel style={/pgf/number format/fixed},
  tick label style={font=\footnotesize},label style={font=\footnotesize},legend style={font=\scriptsize,draw=none,fill=white,fill opacity=0.8,text opacity=1}}}

\title{Bermudan Swaptions in a One-Factor HJM Model\\[4pt]
  \large Calibration, PDE and Monte Carlo Pricing, and a Review of the Model}
\author{Mohammed-Yassine Barnicha}
\date{October 2026}

\begin{document}
\maketitle

\begin{abstract}
\noindent We price co-terminal Bermudan swaptions in the one-factor Heath--Jarrow--Morton model with separable volatility
$\sigma_f(t,T)=\sigma(t)e^{-\lambda(T-t)}$. The model is calibrated exactly to co-terminal European swaptions quoted in normal
volatility, by a sequential bootstrap built on Jamshidian's decomposition, and the Bermudan is priced by a finite-difference
scheme on the Markov factor and by Longstaff--Schwartz Monte Carlo. Reviewing the original implementation, we show that its
zero-coupon bond formula carries the convexity adjustment with the wrong sign: the model then violates the martingale property
of discounted bonds by $1.8\,\%$ on a 7-year bond, prices the same European swaption at $0.0258$ by Jamshidian and $0.0206$ by
the PDE, and undervalues the reference Bermudan by $11\,\%$. We derive the correct formula, measure the other numerical
issues (Monte Carlo discounting, regression bias, time scheme), and study how the mean-reversion parameter, which the
calibration cannot identify, drives the Bermudan price through the correlation between exercise dates. All computations are
reproduced by the interactive page \url{https://yassinebarnicha.github.io/swaptions-hjm/}.
\end{abstract}

\tableofcontents

%======================================================================
\section{Introduction}
A Bermudan swaption gives its holder the right to enter, on any of a set of dates $T_1<\dots<T_m$, an interest-rate swap that
ends at a fixed final date $T_N$. Each exercise right alone is a European swaption, which the market quotes; the Bermudan
adds the right to \emph{choose} the date, and the value of that choice depends on how rates at the different dates are
correlated, which European prices do not reveal. A pricing model must therefore (i) reproduce the Europeans, (ii) have a
realistic dynamics across dates, and (iii) be numerically tractable for an optimal stopping problem.

The one-factor Gaussian HJM model with separable volatility (equivalently the Hull--White model with time-dependent
volatility) meets the third requirement best: it is Markov in one Gaussian factor, bonds are exponential-affine in it,
Europeans have closed forms, and the Bermudan reduces to a one-dimensional PDE. Its mean-reversion speed $\lambda$ sets the
inter-date correlation.

\paragraph{Contributions.} This report documents the project and its review:
\begin{enumerate}[nosep]
  \item a self-contained derivation of the model (\cref{sec:model}), of the European pricing formulas (\cref{sec:europeans})
        and of the calibration (\cref{sec:calibration});
  \item the PDE and Monte Carlo methods for the Bermudan, with their error analysis and convergence study
        (\cref{sec:bermudan,sec:pde,sec:lsmc});
  \item the model review (\cref{sec:review}): a sign error in the bond formula, exposed by a martingale test, and its impact;
  \item a study of mean reversion under recalibration (\cref{sec:meanrev}).
\end{enumerate}

%======================================================================
\section{The model}\label{sec:model}
\subsection{The HJM framework}
Let $P(t,T)$ be the price at $t$ of the zero-coupon bond paying $1$ at $T$, $f(t,T)=-\partial_T\ln P(t,T)$ the instantaneous
forward rate, $r_t=f(t,t)$ the short rate and $\beta_t=\exp\int_0^tr_s\dd s$ the bank account. Under the risk-neutral measure
$\Q$, with a one-dimensional Brownian motion $W$, HJM specify
\begin{equation}
  \dd f(t,T)=\mu(t,T)\dd t+\sigma_f(t,T)\dd W_t,\qquad f(0,T)=f^{\mathrm{mkt}}(0,T).
\end{equation}
\begin{proposition}[HJM drift condition]\label{prop:drift}
Discounted bond prices $P(t,T)/\beta_t$ are $\Q$-martingales for all $T$ if and only if
\begin{equation}\label{eq:drift}
  \mu(t,T)=\sigma_f(t,T)\int_t^T\sigma_f(t,u)\dd u .
\end{equation}
\end{proposition}
\begin{proof}
Let $\Sigma(t,T)=\int_t^T\sigma_f(t,u)\dd u$ and $A(t,T)=\int_t^T\mu(t,u)\dd u$. From $\ln P(t,T)=-\int_t^Tf(t,u)\dd u$ and
the stochastic Fubini theorem, $\dd\ln P(t,T)=\big(r_t-A(t,T)\big)\dd t-\Sigma(t,T)\dd W_t$, so
$\dd\big(P/\beta\big)/(P/\beta)=\big(\tfrac12\Sigma^2-A\big)\dd t-\Sigma\dd W_t$. The drift vanishes for all $T$ iff
$A(t,T)=\tfrac12\Sigma(t,T)^2$; differentiating in $T$ gives \eqref{eq:drift}.
\end{proof}

\subsection{Separable volatility and the Markov factor}
The model uses
\begin{equation}
  \sigma_f(t,T)=\sigma(t)\,e^{-\lambda(T-t)}=\eta(t)\,b(T),\qquad \eta(t)=\sigma(t)e^{\lambda t},\quad b(T)=e^{-\lambda T},
\end{equation}
with $\sigma(\cdot)$ piecewise constant on the exercise grid: $\sigma(t)=\sigma_k$ for $t\in(T_{k-1},T_k]$, $T_0=0$.
Define
\begin{equation}
\begin{gathered}
  X_t=\int_0^t\eta(u)\dd W_u,\qquad
  \alpha(t)=\int_0^t\eta(u)^2\dd u,\\
  \Bb(t)=\int_0^tb(u)\dd u=\frac{1-e^{-\lambda t}}{\lambda},\qquad
  \gamma(t)=\int_0^t\eta(u)^2\,\Bb(u)\dd u .
\end{gathered}
\end{equation}
$X$ is a continuous Gaussian martingale, $X_t\sim\N(0,\alpha(t))$, with independent increments. On a segment $[a,b]$ where
$\sigma\equiv\sigma_j$,
\begin{equation}\label{eq:alphagamma}
  \alpha(b)-\alpha(a)=\sigma_j^2\,\frac{e^{2\lambda b}-e^{2\lambda a}}{2\lambda},\qquad
  \gamma(b)-\gamma(a)=\frac{\sigma_j^2}{\lambda}\left[\frac{e^{2\lambda b}-e^{2\lambda a}}{2\lambda}-\frac{e^{\lambda b}-e^{\lambda a}}{\lambda}\right],
\end{equation}
with the limits $\sigma_j^2(b-a)$ and $\tfrac12\sigma_j^2(b^2-a^2)$ when $\lambda\to0$.

\begin{proposition}[Forward rates, short rate and bonds]\label{prop:bonds}
For $t\le T$,
\begin{align}
  f(t,T)&=f(0,T)+b(T)\big[\Bb(T)\,\alpha(t)-\gamma(t)+X_t\big],\label{eq:fwd}\\
  r_t&=f(0,t)+e^{-\lambda t}\big[\Bb(t)\,\alpha(t)-\gamma(t)\big]+e^{-\lambda t}X_t,\label{eq:short}\\
  P(t,T)&=\frac{P(0,T)}{P(0,t)}\exp\Big(-D(t,T)\big[c(t,T)+X_t\big]\Big),\label{eq:bond}
\end{align}
where $D(t,T)=\Bb(T)-\Bb(t)$ and $c(t,T)=\tfrac12\big(\Bb(T)+\Bb(t)\big)\alpha(t)-\gamma(t)$.
\end{proposition}
\begin{proof}
By \eqref{eq:drift}, $\mu(s,T)=\eta(s)^2b(T)\big(\Bb(T)-\Bb(s)\big)$. Integrating $\dd f(s,T)$ over $s\in[0,t]$,
\[
  f(t,T)=f(0,T)+b(T)\int_0^t\eta(s)^2\big(\Bb(T)-\Bb(s)\big)\dd s+b(T)\int_0^t\eta(s)\dd W_s,
\]
which is \eqref{eq:fwd}; \eqref{eq:short} is $T=t$. For the bond, $\int_t^Tb(u)\dd u=D(t,T)$ and
$\int_t^Tb(u)\Bb(u)\dd u=\tfrac12\big(\Bb(T)^2-\Bb(t)^2\big)=\tfrac12D(t,T)\big(\Bb(T)+\Bb(t)\big)$, so
$\int_t^T\big[f(t,u)-f(0,u)\big]\dd u=D(t,T)\big[c(t,T)+X_t\big]$ and \eqref{eq:bond} follows from
$P(0,T)/P(0,t)=\exp\big(-\int_t^Tf(0,u)\dd u\big)$.
\end{proof}

\begin{remark}[Hull--White]
Applying It\^o's formula to \eqref{eq:short} gives $\dd r_t=\big(\theta(t)-\lambda r_t\big)\dd t+\sigma(t)\dd W_t$ for a
deterministic $\theta(\cdot)$ that fits the initial curve: the model is the Hull--White model with time-dependent volatility.
$\lambda$ is the speed of mean reversion of the short rate and $\sigma(t)$ its normal volatility.
\end{remark}

\begin{remark}[Dynamics of bond ratios]\label{rem:ratio}
From \eqref{eq:bond}, the diffusion part of $\ln P(t,T)$ is $-D(t,T)\,\eta(t)\dd W_t$, hence that of
$\ln\big(P(t,T_i)/P(t,T_e)\big)$ is $-\big(\Bb(T_i)-\Bb(T_e)\big)\eta(t)\dd W_t$: a deterministic constant times $\dd X_t$.
Its integrated variance on $[0,T_e]$ is $D(T_e,T_i)^2\,\alpha(T_e)$.
\end{remark}

%======================================================================
\section{European swaptions}\label{sec:europeans}
\subsection{Swaps and market quotes}
A payer swap starting at $T_0$ with fixed rate $K$, payment dates $T_1<\dots<T_n$ and accruals $\tau_i$ is worth, at $T_0$,
$1-P(T_0,T_n)-K\sum_i\tau_iP(T_0,T_i)$ per unit notional. With the annuity $A_t=\sum_i\tau_iP(t,T_i)$ and the forward
swap rate $S_t=\big(P(t,T_0)-P(t,T_n)\big)/A_t$, the swap value is $A_t(S_t-K)$ and the payer swaption pays
$A_{T_0}(S_{T_0}-K)^+$. Under the annuity measure $\Q^A$ (num\'eraire $A$), $S$ is a martingale. The market convention
(Bachelier) takes $\dd S_t=\sigma_N\dd W^A_t$, so that
\begin{equation}\label{eq:bachelier}
  V^{\mathrm{pay}}_0=A_0\,\E^{A}\big[(S_{T_0}-K)^+\big]=A_0\Big[(S_0-K)\,\Phi(d)+\sigma_N\sqrt{T_0}\,\varphi(d)\Big],\qquad
  d=\frac{S_0-K}{\sigma_N\sqrt{T_0}},
\end{equation}
and $V^{\mathrm{rec}}_0=A_0\big[(K-S_0)\Phi(-d)+\sigma_N\sqrt{T_0}\,\varphi(d)\big]$. A quote $\sigma_N$ is thus a price.

\subsection{Options on zero-coupon bonds}
\begin{proposition}\label{prop:zbo}
The price at $0$ of a put with expiry $T_e$ and strike $K$ on the bond $P(\cdot,T_i)$, $T_i>T_e$, is
\begin{equation}
\begin{gathered}
  \mathrm{ZBP}(T_e,T_i,K)=K\,P(0,T_e)\,\Phi(-h+v)-P(0,T_i)\,\Phi(-h),\\
  v=D(T_e,T_i)\sqrt{\alpha(T_e)},\qquad h=\frac1v\ln\frac{P(0,T_i)}{K\,P(0,T_e)}+\frac v2,
\end{gathered}
\end{equation}
and the call is $\mathrm{ZBC}=P(0,T_i)\Phi(h)-K\,P(0,T_e)\Phi(h-v)$.
\end{proposition}
\begin{proof}
Under the $T_e$-forward measure, $Y_t=P(t,T_i)/P(t,T_e)$ is a martingale; by \cref{rem:ratio} it is lognormal with total
variance $v^2$ at $T_e$, and $Y_{T_e}=P(T_e,T_i)$. Then $\mathrm{ZBP}=P(0,T_e)\,\E^{T_e}[(K-Y_{T_e})^+]$ is Black's formula
with forward $Y_0=P(0,T_i)/P(0,T_e)$.
\end{proof}

\subsection{Jamshidian's decomposition}
\begin{proposition}[Jamshidian]\label{prop:jam}
Let $c_i=K\tau_i$ for $i<n$ and $c_n=1+K\tau_n$. There is a unique $x^\star$ with $\sum_ic_iP(T_0,T_i;x^\star)=1$, and with
$K_i=P(T_0,T_i;x^\star)$,
\begin{equation}
  V^{\mathrm{pay}}_0=\sum_{i=1}^nc_i\,\mathrm{ZBP}(T_0,T_i,K_i),\qquad V^{\mathrm{rec}}_0=\sum_{i=1}^nc_i\,\mathrm{ZBC}(T_0,T_i,K_i).
\end{equation}
\end{proposition}
\begin{proof}
The payer payoff is $\big(1-\sum_ic_iP(T_0,T_i;X_{T_0})\big)^+$. By \eqref{eq:bond}, each $P(T_0,T_i;x)$ is strictly decreasing
in $x$ ($D>0$), from $+\infty$ to $0$, so $\sum_ic_iP(T_0,T_i;x)$ crosses $1$ exactly once, at $x^\star$. For $x\ge x^\star$ every
$P(T_0,T_i;x)\le K_i$ and for $x<x^\star$ every $P(T_0,T_i;x)>K_i$; in both cases
$\big(1-\sum_ic_iP_i\big)^+=\sum_ic_i\big(K_i-P_i\big)^+$, using $\sum_ic_iK_i=1$. Take expectations and apply \cref{prop:zbo}.
\end{proof}
The implementation finds $x^\star$ with Brent's method on $[-10,10]$. Put--call parity (payer $-$ receiver $=$ forward swap)
holds to $10^{-8}$.

%======================================================================
\section{Calibration}\label{sec:calibration}
\subsection{The co-terminal bootstrap}
Each exercise date $T_k$ carries the co-terminal European swaption $T_k\to T_N$, quoted at $\sigma_{N,k}$. By
\cref{prop:jam,prop:zbo}, its model price depends on $\sigma(\cdot)$ only through $\alpha(T_k)$, which by \eqref{eq:alphagamma}
depends on $\sigma_1,\dots,\sigma_k$ only. Each bond put increases with $v$, so the price is continuous and strictly
increasing in $\sigma_k$ once $\sigma_1,\dots,\sigma_{k-1}$ are fixed.
\begin{proposition}\label{prop:boot}
Given $\lambda$ and $\sigma_1,\dots,\sigma_{k-1}$, the equation $V^{\mathrm{model}}_k(\sigma_k)=V^{\mathrm{Bach}}_k(\sigma_{N,k})$
has at most one solution $\sigma_k>0$; it exists iff the market price exceeds the model price with $\sigma_k=0$, i.e.\ iff the
market does not imply less total variance at $T_k$ than is already accumulated on $[0,T_{k-1}]$.
\end{proposition}
The bootstrap solves the dates in increasing order with Brent's method on $[\sigma_{\min},\sigma_{\max}]$, and the fit is
exact (residuals below $10^{-10}$). When the condition of \cref{prop:boot} fails, the page reports the instrument that
cannot be fitted instead of returning a wrong $\sigma$.

\subsection{Mean reversion is not identified}
For every $\lambda$ in a wide range, \cref{prop:boot} produces a $\sigma(\cdot)$ that reproduces all the co-terminal Europeans
exactly. The calibration set therefore leaves a one-parameter family of perfectly calibrated models, whose Bermudan prices
differ (\cref{sec:meanrev}): on our reference case, from $0.01656$ ($\lambda=0$) to $0.01848$ ($\lambda=0.3$), a spread of
$12\,\%$. $\lambda$ must be set from other information: historical correlation of rates, or other instruments (caps, diagonal or
column swaptions) whose prices depend on inter-date correlation.

%======================================================================
\section{The Bermudan swaption}\label{sec:bermudan}
The exercise value at $T_k$ is the underlying co-terminal swap, $h_k(x)=\big(1-P(T_k,T_N;x)-K\sum_{T_i>T_k}\tau_iP(T_k,T_i;x)\big)^+$
for a payer. The price is the value of an optimal stopping problem over the stopping times $\tau$ taking values in
$\{T_1,\dots,T_m\}$:
\begin{equation}
  V_0=\sup_\tau\E^\Q\big[\beta_\tau^{-1}h_\tau(X_\tau)\big].
\end{equation}
Since $X$ is Markov, dynamic programming gives
\begin{equation}\label{eq:dp}
  V_{T_m}=h_m,\qquad V_{T_k}=\max\big(h_k,C_k\big),\qquad
  C_k(x)=\E^\Q\Big[e^{-\int_{T_k}^{T_{k+1}}r_s\dd s}\,V_{T_{k+1}}\,\Big|\,X_{T_k}=x\Big].
\end{equation}
Exercising always at a fixed $T_k$ is admissible, so $V_0\ge\max_kV^{\mathrm{Eur}}_k$; the difference is the
\emph{early-exercise} (or switch) premium. On the reference case below the Bermudan is worth $0.016746$ against
$0.014328$ for the most expensive European.

\paragraph{Reference case.} Curve $P(0,t)$ at $t=1,2,3,5,7,10$: $0.978$, $0.954$, $0.930$, $0.885$, $0.842$, $0.780$; payer, $K=3\,\%$,
annual payments from $2$ to $T_N=7$; exercise dates $1,2,3$ with normal vols $80$, $85$, $90$\,bp; $\lambda=3\,\%$. The
bootstrap gives $\sigma=(86.12,\ 96.45,\ 106.37)$\,bp.

%======================================================================
\section{The PDE method}\label{sec:pde}
\subsection{Equation and boundary conditions}
Between exercise dates, by Feynman--Kac, $V(t,x)=\E[e^{-\int_t^{T}r}V(T,X_T)\,|\,X_t=x]$ solves
\begin{equation}\label{eq:pde}
  \partial_tV+\tfrac12\eta(t)^2\,\partial_{xx}V-r(t,x)\,V=0,\qquad r(t,x)=A(t)+e^{-\lambda t}x,
\end{equation}
with $A(t)=f(0,t)+e^{-\lambda t}\big[\Bb(t)\alpha(t)-\gamma(t)\big]$ from \eqref{eq:short}. There is no drift term because $X$ is
a martingale; this is why the factor $X$, rather than $r$, is the natural state variable. At each $T_k$ the value is
updated by $V\leftarrow\max(V,h_k)$, and the equation is solved backwards from $V(T_m,\cdot)=h_m$.

The domain is truncated to $[-n\sqrt{\alpha(T_m)},\,n\sqrt{\alpha(T_m)}]$. Far from the origin the option is either deep in
the money (exercise immediately) or worthless, so Dirichlet values equal to the largest intrinsic value among the remaining
exercise dates are imposed. The truncation error is of the order of the probability of leaving the domain,
$2\Phi(-n)\approx e^{-n^2/2}$, which is negligible from $n=5$ (\cref{tab:domain}).

\subsection{Discretisation}
On a uniform grid $x_j$ with step $\Delta x$, let $\mathcal L^n$ be the discrete operator
$(\mathcal L^nV)_j=\tfrac12\eta(t_n)^2\,(V_{j+1}-2V_j+V_{j-1})/\Delta x^2-r(t_n,x_j)V_j$. The $\theta$-scheme reads
\begin{equation}
  \frac{V^n-V^{n+1}}{\Delta t}=\theta\,\mathcal L^nV^n+(1-\theta)\,\mathcal L^{n+1}V^{n+1},
\end{equation}
a tridiagonal system per step, solved by the Thomas algorithm in $O(n_x)$ operations. The time grid contains every exercise
date, so that the $\max$ is applied exactly.
\begin{itemize}[nosep]
  \item $\theta=1$ (implicit Euler, the original code) is unconditionally stable and monotone whenever
        $1+\Delta t\,r(t,x)>0$, but only first order in time.
  \item $\theta=\tfrac12$ (Crank--Nicolson) is second order and unconditionally stable in $L^2$, but not L-stable: the
        high-frequency components created by the kink of $h_k$ at each exercise date are damped by a factor close to $-1$ per
        step and oscillate.
  \item \emph{Rannacher time-stepping} replaces the first two Crank--Nicolson steps after each exercise date by implicit
        steps (here: two full steps). This damps the oscillations and restores second-order convergence for non-smooth data
        (Rannacher 1984; Giles and Carter 2006).
\end{itemize}

\subsection{Convergence}
\Cref{fig:conv} shows the error in time against a reference computed with $n_t=12\,800$, and the error in space. Implicit
Euler converges at order one and Crank--Nicolson with Rannacher steps at order two (the error is divided by about $4$ to $5$
each time $n_t$ doubles). On this price, plain Crank--Nicolson is even more accurate at coarse steps: the oscillations it
leaves near the exercise boundary are small and partly cancel in the value at $x=0$; they matter more for the Greeks, which is
why the Rannacher steps are kept by default. In space the centred scheme is second order. With $n_x=501$, $n_t=1000$ and
Crank--Nicolson--Rannacher the price is accurate to a few $10^{-7}$ (space error) and costs a few tens of milliseconds in the
browser.

This study also caught an error in an early version of the port: the $\theta$-scheme used $\eta(t_n)$ in both halves of
the step, a first-order error in time since $\eta(t)=\sigma(t)e^{\lambda t}$ varies within a step. The scheme now uses
$\eta(t_n)$ in the implicit part and $\eta(t_{n+1})$ in the explicit part, with $\sigma$ read on the segment of the step.

\begin{figure}[ht]\centering
\begin{tikzpicture}
\begin{loglogaxis}[xlabel={time steps $n_t$},ylabel={$|V-V_{\mathrm{ref}}|$},legend pos=north east]
  \addplot+[mark=*,thick] table[x=nt,y=implicite]{data/temps.dat}; \addlegendentry{implicit}
  \addplot+[mark=square*,thick] table[x=nt,y=cn]{data/temps.dat}; \addlegendentry{Crank--Nicolson}
  \addplot+[mark=triangle*,thick] table[x=nt,y=cnr]{data/temps.dat}; \addlegendentry{CN + Rannacher}
\end{loglogaxis}
\end{tikzpicture}\hfill
\begin{tikzpicture}
\begin{loglogaxis}[xlabel={space points $n_x$},ylabel={$|V-V_{\mathrm{ref}}|$}]
  \addplot+[mark=*,thick] table[x=nx,y=err]{data/espace.dat};
\end{loglogaxis}
\end{tikzpicture}
\caption{PDE convergence on the reference case. Left: time error ($n_x=1001$). Right: space error (CN + Rannacher,
$n_t=4000$, reference $n_x=4001$).}\label{fig:conv}
\end{figure}

\begin{table}[ht]\centering\small
\pgfplotstabletypeset[columns={nstd,prix},columns/nstd/.style={column name={$n$ (standard deviations)}},
  columns/prix/.style={column name={Bermudan price},fixed,precision=8},every head row/.style={before row=\toprule,after row=\midrule},
  every last row/.style={after row=\bottomrule}]{data/domaine.dat}
\caption{Truncation of the PDE domain to $\pm n\sqrt{\alpha(T_m)}$ (space step kept constant).}\label{tab:domain}
\end{table}

\subsection{Exercise boundary}
At each date the holder exercises where $h_k(x)>C_k(x)$. For a payer this is a half-line $x>x^\star_k$: high factor, high
rates, valuable payer swap. \Cref{fig:boundary} shows $C_k$ and $h_k$ at $T_1$ and $T_2$; at the last date $C_m=0$ and the
option is exercised as soon as it is in the money.

\begin{figure}[ht]\centering
\begin{tikzpicture}
\begin{axis}[xlabel={factor $x$},ylabel={value},title={\footnotesize $T_1=1$},legend pos=north west,ymin=-0.002,xmin=-0.06,xmax=0.075]
  \addplot[blue,thick] table[x=x,y=c1]{data/frontiere.dat}; \addlegendentry{continuation $C_1$}
  \addplot[red,thick,dashed] table[x=x,y=e1]{data/frontiere.dat}; \addlegendentry{exercise $h_1$}
\end{axis}
\end{tikzpicture}\hfill
\begin{tikzpicture}
\begin{axis}[xlabel={factor $x$},title={\footnotesize $T_2=2$},ymin=-0.002,xmin=-0.06,xmax=0.075]
  \addplot[blue,thick] table[x=x,y=c2]{data/frontiere.dat};
  \addplot[red,thick,dashed] table[x=x,y=e2]{data/frontiere.dat};
\end{axis}
\end{tikzpicture}
\caption{Continuation and exercise values from the PDE. The exercise boundary is the crossing point.}\label{fig:boundary}
\end{figure}

%======================================================================
\section{Longstaff--Schwartz Monte Carlo}\label{sec:lsmc}
\subsection{Algorithm}
The factor is simulated exactly: $X_{t+\Delta}-X_t\sim\N\big(0,\alpha(t+\Delta)-\alpha(t)\big)$, independent of the past. The
continuation value in \eqref{eq:dp} is a conditional expectation, estimated by least squares (Longstaff and Schwartz 2001).

\begin{algorithm}[ht]\caption{Longstaff--Schwartz for the Bermudan swaption}\label{alg:lsmc}
\begin{algorithmic}[1]
\State simulate $N$ paths of $X$ (and of $I_t=\int_0^tr_s\dd s$) at the exercise dates
\State $\mathrm{CF}^{(i)}\gets h_m(X^{(i)}_{T_m})$, $\ \tau^{(i)}\gets T_m$ for every path $i$
\For{$k=m-1,\dots,1$}
  \State $\mathcal I\gets\{i: h_k(X^{(i)}_{T_k})>0\}$ \Comment{in-the-money paths only}
  \State regress $Y^{(i)}=e^{-(I^{(i)}_{\tau^{(i)}}-I^{(i)}_{T_k})}\,\mathrm{CF}^{(i)}$ on $\{1,\tilde x,\dots,\tilde x^d\}$,
        $\tilde x=x/\sqrt{\alpha(T_k)}$, over $i\in\mathcal I$, giving $\hat C_k$
  \For{$i\in\mathcal I$ with $h_k(X^{(i)}_{T_k})>\hat C_k(X^{(i)}_{T_k})$}
     \State $\mathrm{CF}^{(i)}\gets h_k(X^{(i)}_{T_k})$, $\ \tau^{(i)}\gets T_k$
  \EndFor
\EndFor
\State \Return $\hat V_0=\frac1N\sum_ie^{-I^{(i)}_{\tau^{(i)}}}\,\mathrm{CF}^{(i)}$ and its standard error
\end{algorithmic}
\end{algorithm}

Normalising the regressor by $\sqrt{\alpha(T_k)}$ does not change the fit but keeps the normal equations well conditioned;
the least-squares problem is solved by a Householder QR factorisation.

\subsection{The two biases}
Write $\hat\tau$ for the exercise rule produced by the regression.
\begin{itemize}[nosep]
  \item \emph{Low bias.} For any stopping rule, $\E[\beta_{\hat\tau}^{-1}h_{\hat\tau}]\le V_0$. If $\hat\tau$ is estimated on
        one set of paths and applied to an independent set, the second set gives an unbiased estimate of this quantity, hence a
        \emph{lower bound} of $V_0$ in expectation.
  \item \emph{High bias.} When the same paths are used to estimate $\hat\tau$ and to price, the rule is fitted to the noise of
        those paths and anticipates their future, which pushes the estimate up.
\end{itemize}
The page offers both estimators. An upper bound can be obtained by duality (Rogers 2002; Haugh and Kogan 2004; Andersen and
Broadie 2004), a natural extension.

\subsection{Discounting}
Discounting must use the path of the short rate: $\beta_{T_k}^{-1}=e^{-I_{T_k}}$. The implementation simulates $X$ on a fine
grid ($52$ steps per year by default) and integrates $r(t,X_t)$ by the trapezoidal rule, an $O(\Delta t^2)$ error. An exact
alternative exists: $(X_t,I_t)$ is a Gaussian vector whose mean and covariance follow from \eqref{eq:short}, so it could be
sampled jointly at the exercise dates without any grid.

\subsection{Results}
\Cref{fig:lsmc} shows the three estimators against the PDE price. With exact discounting, both the in-sample and the
independent estimators contain the PDE value in their confidence intervals: with a degree-2 regression on one factor, both
biases are smaller than the statistical error at these sample sizes. The original discounting (\cref{sec:review}) stays
about $3$\,bp of notional above the PDE whatever the number of paths: it is a bias, not noise.

\begin{figure}[ht]\centering
\begin{tikzpicture}
\begin{semilogxaxis}[width=0.75\textwidth,height=6cm,xlabel={paths $N$},ylabel={Bermudan price},legend pos=north east,
  yticklabel style={/pgf/number format/fixed,/pgf/number format/precision=3},xticklabel style={/pgf/number format/sci}]
  \addplot+[mark=*,thick,error bars/.cd,y dir=both,y explicit] table[x=N,y=dans,y error=dans_ic]{data/lsmc.dat}; \addlegendentry{exact discounting, in-sample}
  \addplot+[mark=square*,thick,error bars/.cd,y dir=both,y explicit] table[x=N,y=indep,y error=indep_ic]{data/lsmc.dat}; \addlegendentry{exact discounting, independent paths}
  \addplot+[mark=triangle*,thick,error bars/.cd,y dir=both,y explicit] table[x=N,y=code,y error=code_ic]{data/lsmc.dat}; \addlegendentry{original discounting}
  \addplot[black,dashed,domain=1000:100000] {0.016746}; \addlegendentry{PDE}
\end{semilogxaxis}
\end{tikzpicture}
\caption{Longstaff--Schwartz estimates with $95\,\%$ confidence intervals, degree-2 regression, against the PDE price.}\label{fig:lsmc}
\end{figure}

%======================================================================
\section{Review of the original model}\label{sec:review}
\subsection{A sign error in the convexity term}
The original \texttt{hjm\_1f.py} implements
\begin{equation}
  P^{\mathrm{code}}(t,T)=\frac{P(0,T)}{P(0,t)}\exp\Big(D(t,T)\big[c(t,T)-X_t\big]\Big),
\end{equation}
that is \eqref{eq:bond} with $+c$ instead of $-c$, and its short rate, $r^{\mathrm{code}}_t=-\partial_T\ln P^{\mathrm{code}}(t,T)|_{T=t}$,
inherits the error: $r^{\mathrm{code}}_t=f(0,t)-e^{-\lambda t}\big[\Bb(t)\alpha(t)-\gamma(t)\big]+e^{-\lambda t}X_t$.
\Cref{prop:bonds} derives the correct sign directly from the drift condition. Because $c(t,T)>0$ grows with
$\alpha(t)$, the error grows with volatility and maturity.

\paragraph{Why calibration does not see it.} Jamshidian's formula (\cref{prop:jam}) uses $P(0,\cdot)$ and the variance
$v^2$ of bond ratios, neither of which depends on $c$: Europeans are priced correctly and the calibrated $\sigma$ are the
same to $10^{-2}$\,bp. The error only appears in methods that use $r(t,x)$ and $P(t,T;x)$ along paths, the PDE and the Monte
Carlo, so the model prices the same European differently depending on the method.

\subsection{The martingale test}
Under $\Q$, $\E\big[\beta_t^{-1}P(t,T)\big]$ must equal $P(0,T)$ for every $t$. \Cref{fig:mart} estimates it for $T=7$ with
$200\,000$ paths ($\sigma=150$\,bp, $\lambda=3\,\%$). The original formula drifts away, to $2.2\,\%$ at $t=6$; the corrected
one stays within the confidence band.

\begin{figure}[ht]\centering
\begin{tikzpicture}
\begin{axis}[width=0.72\textwidth,height=5.6cm,xlabel={$t$ (years)},ylabel={$\E[\beta_t^{-1}P(t,7)]/P(0,7)-1$ (\%)},legend pos=north west]
  \addplot+[mark=*,thick,error bars/.cd,y dir=both,y explicit] table[x=t,y=err_pct,y error=ic_pct]{data/martingale_code.dat};
  \addlegendentry{original formula}
  \addplot+[mark=square*,thick,error bars/.cd,y dir=both,y explicit] table[x=t,y=err_pct,y error=ic_pct]{data/martingale_corr.dat};
  \addlegendentry{corrected formula}
\end{axis}
\end{tikzpicture}
\caption{Martingale test of discounted bond prices, with $95\,\%$ confidence intervals.}\label{fig:mart}
\end{figure}

\subsection{Impact on prices}
\begin{table}[ht]\centering\small
\begin{tabular}{@{}lll@{}}\toprule
 & original formula & corrected formula\\\midrule
$\E[\beta_3^{-1}P(3,7)]$, vs $P(0,7)=0.842000$ & $0.856942$ ($+1.78\,\%$, $83$ s.e.) & $0.841949$ ($-0.3$ s.e.)\\
European $3\mathrm{y}\to7\mathrm{y}$, Jamshidian & $0.02582$ & $0.02582$\\
European $3\mathrm{y}\to7\mathrm{y}$, PDE & $0.02057$ & $0.02582$\\
European $3\mathrm{y}\to7\mathrm{y}$, Monte Carlo & $0.02049\pm0.00015$ & $0.02574\pm0.00017$\\
Reference Bermudan (PDE, implicit) & $0.014929$ & $0.016737$\\
Early-exercise premium & $0.000601$ & $0.002409$\\\bottomrule
\end{tabular}
\caption{Consequences of the sign error. Rows 1--4: $\sigma=150$\,bp, $\lambda=3\,\%$. Rows 5--6: reference case
($n_x=501$, $n_t=1000$); the premium is the Bermudan minus the most expensive European, $0.014328$, which is the same with both
formulas.}\label{tab:impact}
\end{table}

The last row is the most telling. Since Europeans are unaffected, the whole error falls on the early-exercise premium, the
quantity a Bermudan model exists to price: the original formula underestimates it by $75\,\%$.

\subsection{Other findings}
\begin{itemize}
  \item \emph{Monte Carlo discounting.} \texttt{mc\_1f.py} freezes the short rate between exercise dates and discounts the
        exercise value back to today at $r(0)$, ignoring the curve: on the reference case, Monte Carlo bond prices are off by
        $0.26\,\%$ at $2$ years and $0.59\,\%$ at $3$ years, and the Bermudan is overpriced by about $3$\,bp of notional.
        Fixed by the exact discounting of \cref{sec:lsmc}.
  \item \emph{Regression bias.} Pricing on the regression paths only; independent paths and a choice of degree were added.
  \item \emph{Time scheme.} Implicit Euler only; Crank--Nicolson with Rannacher steps was added (\cref{fig:conv}).
  \item \emph{Receiver swaptions.} The Python returns the payer formula for a receiver; the port uses bond calls.
\end{itemize}

%======================================================================
\section{Mean reversion under recalibration}\label{sec:meanrev}
\subsection{Correlation between exercise dates}
Since $X$ has independent increments, $\operatorname{Cov}(X_s,X_t)=\alpha(s)$ for $s\le t$ and
\begin{equation}\label{eq:corr}
  \rho(s,t)=\operatorname{corr}(X_s,X_t)=\sqrt{\alpha(s)/\alpha(t)} .
\end{equation}
The short rate and, to first order, the swap rates are affine in $X$ with deterministic coefficients, so \eqref{eq:corr} is
also their correlation across dates.
\begin{proposition}\label{prop:corr}
For a constant $\sigma$ and $0<s<t$, $\rho(s,t)^2=\dfrac{e^{2\lambda s}-1}{e^{2\lambda t}-1}$ is strictly decreasing in
$\lambda>0$, from $s/t$ at $\lambda=0$ to $0$ as $\lambda\to\infty$.
\end{proposition}
\begin{proof}
$\partial_\lambda\ln\rho^2=g(s)-g(t)$ with $g(u)=\dfrac{2u\,e^{2\lambda u}}{e^{2\lambda u}-1}=\dfrac1\lambda\,\psi(2\lambda u)$
and $\psi(z)=z/(1-e^{-z})$, which is strictly increasing on $(0,\infty)$. Hence $g(s)<g(t)$. The limits follow from
$e^{2\lambda u}-1\sim2\lambda u$ and from $\rho^2\le e^{-2\lambda(t-s)}\cdot\frac{1}{1-e^{-2\lambda t}}$.
\end{proof}

\subsection{Why the Bermudan becomes more valuable}
Recalibration pins each European: in the swap-rate approximation, the $k$-th European fixes the model variance of its own
swap rate, $\alpha(T_k)\,\kappa_k(\lambda)^2\approx\sigma_{N,k}^2T_k$, where $\kappa_k=\partial S_k/\partial x$ is a duration-like
factor containing $e^{-\lambda T_k}$. A higher $\lambda$ damps the impact of past shocks, so the calibration raises
$\sigma$ (\cref{fig:mr}, left: from $(77.9,\,87.3,\,96.7)$\,bp at $\lambda=0$ to $(188.9,\,214.4,\,228.9)$\,bp at $\lambda=0.3$).
The marginal law of each exercise value is then essentially unchanged, but the dates decorrelate:
$\rho(T_1,T_3)$ falls from $0.513$ to $0.349$. With the marginals fixed and the dependence weaker, the right to choose the
best date is worth more, exactly as the value of an option on the maximum of two assets increases when their correlation
decreases. The Bermudan rises from $0.01656$ to $0.01848$ while every European stays at its market price
(\cref{fig:mr}, right).

Without recalibration ($\sigma$ frozen at the $\lambda=3\,\%$ values) the effect reverses: a higher $\lambda$ only reduces the
variance of every future rate, and every price falls.

\begin{figure}[ht]\centering
\begin{tikzpicture}
\begin{axis}[xlabel={$\lambda$},ylabel={calibrated $\sigma_k$ (bp)},legend pos=north west]
  \addplot+[thick,mark=none] table[x=lam,y=s1]{data/lambda.dat}; \addlegendentry{$\sigma_1$}
  \addplot+[thick,mark=none] table[x=lam,y=s2]{data/lambda.dat}; \addlegendentry{$\sigma_2$}
  \addplot+[thick,mark=none] table[x=lam,y=s3]{data/lambda.dat}; \addlegendentry{$\sigma_3$}
\end{axis}
\end{tikzpicture}\hfill
\begin{tikzpicture}
\begin{axis}[xlabel={$\lambda$},ylabel={price},legend columns=2,legend style={at={(0.5,-0.3)},anchor=north},
  yticklabel style={/pgf/number format/fixed,/pgf/number format/precision=3}]
  \addplot[blue,thick] table[x=lam,y=berm]{data/lambda.dat}; \addlegendentry{Berm., recalibrated}
  \addplot[blue,thick,dashed] table[x=lam,y=bermfixe]{data/lambda.dat}; \addlegendentry{Berm., $\sigma$ frozen}
  \addplot[red,thick] table[x=lam,y=eur]{data/lambda.dat}; \addlegendentry{max Eur., recalibrated}
  \addplot[red,thick,dashed] table[x=lam,y=eurfixe]{data/lambda.dat}; \addlegendentry{max Eur., $\sigma$ frozen}
\end{axis}
\end{tikzpicture}

\begin{tikzpicture}
\begin{axis}[xlabel={$\lambda$},ylabel={correlation},legend pos=north east]
  \addplot+[thick,mark=none] table[x=lam,y=corr12]{data/lambda.dat}; \addlegendentry{$\rho(T_1,T_2)$}
  \addplot+[thick,mark=none] table[x=lam,y=corr13]{data/lambda.dat}; \addlegendentry{$\rho(T_1,T_3)$}
\end{axis}
\end{tikzpicture}\hfill
\begin{tikzpicture}
\begin{axis}[xlabel={$\rho(T_1,T_3)$},ylabel={Bermudan $-$ max European},x dir=reverse,
  yticklabel style={/pgf/number format/fixed,/pgf/number format/precision=4},xticklabel style={/pgf/number format/precision=2}]
  \addplot+[thick,mark=*,mark size=1.2pt] table[x=corr13,y expr=\thisrow{berm}-\thisrow{eur}]{data/lambda.dat};
\end{axis}
\end{tikzpicture}
\caption{Mean reversion with recalibration. Top: calibrated volatilities and prices. Bottom: correlation between exercise
dates and the early-exercise premium against it ($\lambda$ from $0$ to $0.3$).}\label{fig:mr}
\end{figure}

\begin{remark}
This is a tendency, not a theorem: it relies on the calibration set (co-terminal Europeans), on the payoff, and on the
one-factor structure. In a two-factor model, correlation between dates can be lowered without changing $\lambda$, which is one
reason why two-factor models are preferred for products sensitive to inter-date correlation.
\end{remark}

%======================================================================
\section{Validation}
The JavaScript engine of the page is a line-by-line port of the Python package. With the original formula and the notebook
settings (implicit scheme, $n_x=501$, $n_t=1000$, $\pm5$ standard deviations) it reproduces every PDE price of the notebook
to the eighth decimal, for all grids, domains, $\lambda\in[0,0.3]$ and $K\in[2\,\%,4\,\%]$, and the calibrated $\sigma$ to
$10^{-17}$. With the corrected formula, Jamshidian, the PDE and the Monte Carlo agree on European swaptions, and simulated bond
prices recover the curve within their confidence intervals.

\section{Conclusion}
The one-factor HJM model calibrates exactly and quickly to co-terminal Europeans and prices Bermudans accurately by PDE, with
Monte Carlo as an independent check. The review shows the value of model-level tests: a convexity sign error invisible to the
calibration was caught by a martingale test and by the inconsistency between methods, and it moved the Bermudan price by
$12\,\%$. The remaining model risk sits in $\lambda$, which the calibration does not determine and which moves the Bermudan by
$12\,\%$ over a plausible range. Natural extensions are an exact joint simulation of $(X,\int r)$, a dual upper bound for the
Monte Carlo, the calibration of $\lambda$ to correlation-sensitive instruments, and a two-factor version.

\begin{thebibliography}{99}\small
\bibitem{and} L.~Andersen and M.~Broadie. Primal-dual simulation algorithm for pricing multidimensional American options.
  \emph{Management Science}, 50(9):1222--1234, 2004.
\bibitem{ap} L.~Andersen and V.~Piterbarg. \emph{Interest Rate Modeling}, vols.~I--III. Atlantic Financial Press, 2010.
\bibitem{bm} D.~Brigo and F.~Mercurio. \emph{Interest Rate Models -- Theory and Practice}. Springer, 2nd ed., 2006.
\bibitem{gc} M.~Giles and R.~Carter. Convergence analysis of Crank--Nicolson and Rannacher time-marching.
  \emph{Journal of Computational Finance}, 9(4):89--112, 2006.
\bibitem{hk} M.~Haugh and L.~Kogan. Pricing American options: a duality approach. \emph{Operations Research}, 52(2):258--270, 2004.
\bibitem{hjm} D.~Heath, R.~Jarrow and A.~Morton. Bond pricing and the term structure of interest rates: a new methodology for
  contingent claims valuation. \emph{Econometrica}, 60(1):77--105, 1992.
\bibitem{hw} J.~Hull and A.~White. Pricing interest-rate-derivative securities. \emph{Review of Financial Studies}, 3(4):573--592, 1990.
\bibitem{jam} F.~Jamshidian. An exact bond option formula. \emph{Journal of Finance}, 44(1):205--209, 1989.
\bibitem{ls} F.~Longstaff and E.~Schwartz. Valuing American options by simulation: a simple least-squares approach.
  \emph{Review of Financial Studies}, 14(1):113--147, 2001.
\bibitem{ran} R.~Rannacher. Finite element solution of diffusion problems with irregular data. \emph{Numerische Mathematik},
  43:309--327, 1984.
\bibitem{rs} P.~Ritchken and L.~Sankarasubramanian. Volatility structures of forward rates and the dynamics of the term
  structure. \emph{Mathematical Finance}, 5(1):55--72, 1995.
\bibitem{rog} L.~C.~G.~Rogers. Monte Carlo valuation of American options. \emph{Mathematical Finance}, 12(3):271--286, 2002.
\end{thebibliography}
\end{document}
